\pdfoutput=1
\documentclass[11pt]{amsart}
\usepackage{amsmath,amssymb,amsthm,stmaryrd,mathrsfs}
\usepackage[margin=1.2in]{geometry}
\usepackage{xcolor,xurl}
\usepackage[colorlinks=true,linkcolor=blue,citecolor=blue,urlcolor=blue]{hyperref}
\hypersetup{pdfauthor={Bernd Johannes Wuebben},
  pdftitle={Antidiagonal local curves and integral exponents in universal curve counting}}
\setlength{\emergencystretch}{2em}
\newtheorem{theorem}{Theorem}[section]
\newtheorem{proposition}[theorem]{Proposition}
\newtheorem{lemma}[theorem]{Lemma}
\newtheorem{corollary}[theorem]{Corollary}
\theoremstyle{remark}
\newtheorem{remark}[theorem]{Remark}
\newcommand{\ZZ}{\mathbb Z}
\newcommand{\QQ}{\mathbb Q}
\newcommand{\CC}{\mathbb C}
\newcommand{\PP}{\mathbb P}
\newcommand{\Z}{\mathcal Z}
\newcommand{\cO}{\mathcal O}
\DeclareMathOperator{\GW}{GW}
\DeclareMathOperator{\PT}{PT}
\DeclareMathOperator{\Tot}{Tot}
\title[Antidiagonal local curves and universal counting]{Antidiagonal local curves and integral exponents in universal curve counting}
\author{Bernd Johannes Wuebben}
\subjclass[2020]{14N35 (primary); 14J32, 53D45 (secondary)}
\keywords{Universal curve counting, Calabi--Yau threefolds, local curves,
equivariant Gromov--Witten theory, stable pairs, integrality}
\date{November 11, 2026\\This version: November 22, 2026}
\begin{document}
\begin{abstract}
Assuming Pardon's Gromov--Witten evaluation for analytic families of
threefolds, we lift the elementary-cluster expansion to his rational
algebra of one-cycles. For each genus we construct
a universal local factor from the circle action with opposite weights
on the two summands of a balanced local Calabi--Yau curve.
For every smooth projective Calabi--Yau threefold with torsion-free
second homology, the universal curve-counting series is a unique product
of these factors with integer exponents. The exponents are the
elementary-cluster coefficients of Ionel and Parker.
To construct the factors, we project a cycle onto the two line summands
and use the resulting symmetric functions to isolate each multiple of
the zero section. The associated compactly supported classes respect
addition of cycles and multiplication by positive integers, and their Gromov--Witten
evaluations are the local series of Bryan and Pandharipande.
The corresponding stable-pair evaluation gives a second proof of
exponent integrality.
\end{abstract}

\maketitle
\section{Introduction}

Let $C$ be a smooth projective curve of genus $g$ and let $Y\to C$ be
a rank-two vector bundle with determinant $K_C$. Two circle actions
on the total space give different answers to a question about
integrality. Scalar multiplication is intrinsic to the bundle.
For a splitting into two lines the fiber-scaling torus acts on a
holomorphic volume form by pullback with character $(s,t)\mapsto st$.
Its kernel is the antidiagonal torus. Up to inversion its only faithful
one-parameter parametrization has weights $(1,-1)$. For a splitting of degrees
$(g-1,g-1)$, called \emph{balanced}, Bryan and Pandharipande computed
the antidiagonal Gromov--Witten series \cite[Corollary~7.3]{BP}.
The same opposite weights occur in the normal operator of an elementary
cluster. In the spin model of Ionel and Parker the normal operator
is the direct sum of two copies of a real Cauchy--Riemann operator.
Real rotations of these copies commute with that operator. In the
fiber coordinates $v_1+i v_2$ and $v_1-i v_2$ the rotations have
characters $t$ and $t^{-1}$. In the obstruction-bundle calculation,
specializing to these weights gives the Chern class signs of the dual
index bundle and hence the elementary-cluster contribution
\cite[Propositions~3.2--3.3, (3.2)--(3.3)]{IP}.

Pardon associates a universal class to each compact open-and-closed
locus in the space of one-cycles on a complex threefold.
Gromov--Witten and stable-pair invariants
are evaluations of this class. His scalar equivariant local classes
freely generate the Chern-zero, virtual-dimension-zero part of the
rational semi-Fano algebra \cite[Theorem~1.2]{Pardon}. For the local
series formed from those generators, Conjecture~1.4 of \cite{Pardon}
predicts a product expansion with integer exponents. The compact
example of \cite{Counterexample} has a scalar exponent in
$\ZZ+2/3$. It leaves open whether the elementary local series
themselves have geometric representatives in the universal algebra.
We construct such representatives using the antidiagonal action.

Write $H$ for the Chern-zero, virtual-dimension-zero part of
Pardon's rational semi-Fano algebra. For a smooth projective
Calabi--Yau threefold $X$ over $\CC$, let $\Gamma_X$ be its monoid of
effective integral curve classes. We use the convention
$K_X\simeq\cO_X$ and $H^1(X,\cO_X)=H^2(X,\cO_X)=0$.
The coefficient of $Q^\beta$ in the universal series $\Z_X$ is the
class of the characteristic function of the compact cycle locus of
class $\beta$; the constant term is $1$.
We complete the monoid ring in degree against an ample divisor.
Write $c_{\beta,g}(X)$ for the coefficient of the genus-$g$
elementary cluster at class $\beta$ in the connected Gromov--Witten
expansion \cite[Theorem~8.1]{IP}. It multiplies the logarithm of
the genus-$g$ local hook series \eqref{eq:hook} evaluated at $Q^\beta$.

For the identification of the local factors and the compact product
formula we assume Pardon's universal Gromov--Witten evaluation
\cite[\S\S3.4--3.5]{Pardon}. Its definition requires stable-map spaces
over arbitrary separated submersions of reduced complex analytic
spaces of relative dimension three. Pardon indicates how the
construction for algebraic families should extend and notes that he
knows no reference for this step \cite[Remark~1.8]{Pardon}.

\begin{theorem}\label{thm:main}
Assume the universal Gromov--Witten evaluation for analytic families.
Let $C$ be a smooth connected projective curve of genus $g$ and
let $L$ be a line bundle of degree $g-1$ on $C$.
The antidiagonal action on $L\oplus(K_C\otimes L^{-1})$ defines
classes $a_{g,m}\in H$ for $m\geq 1$ which depend only on $g$ and
$m$. Put $A_g(q)=1+\sum_{m\geq1}a_{g,m}q^m$.
For every smooth projective Calabi--Yau threefold $X$,
\begin{equation}\label{eq:main}
 \Z_X=\prod_{\beta\in\Gamma_X\setminus\{0\}}\prod_{g\geq0}
             A_g(Q^\beta)^{c_{\beta,g}(X)} .
\end{equation}
The exponents are unique, integral, and zero for all sufficiently
large $g$ at each fixed $\beta$.
\end{theorem}

The integer exponents and their genus finiteness are known for the
\emph{numerical} elementary expansion
\cite[Theorem~8.1]{IP}, \cite[Theorem~3.50]{DIW}. The assertion of
\eqref{eq:main} is the identity in $H\llbracket\Gamma_X\rrbracket$.
The construction must therefore respect the operations on cycles
before numerical evaluation. For a cycle of degree $m$ over $C$,
the elementary symmetric functions of its two fiber projections
have opposite weights. Their common vanishing over sufficiently
many points of $C$ isolates $m[C]$. The relative Euler classes of
these functions can be normalized for either sign of the weight.
In the coproduct calculation the two summands of a cycle lie over
independent Borel bases. The diagonal of those bases has just one
K\"unneth summand that survives in virtual dimension zero.
Reconstruction identifies the resulting factor as the unique series
compatible with both operations whose evaluation is the elementary
cluster series (Corollary~\ref{cor:characterization}). For compact $X$
we keep the full class in $H_2(X;\ZZ)$. When several effective classes
have the same $d$-fold multiple, pullback under multiplication by $d$
sums their contributions.

Sections~\ref{sec:local} and~\ref{sec:operations} construct the classes
and prove their compatibility with cycle operations independently of the
family evaluations. The second integrality argument in
Section~\ref{sec:factor} assumes the corresponding stable-pair evaluation.

\section{The antidiagonal local classes}\label{sec:local}

Set $L_+=L$, $L_-=K_C\otimes L^{-1}$ and
$Y=\Tot(L_+\oplus L_-)$. The torus acts trivially on $C$ and by
$(v_+,v_-)\mapsto(tv_+,t^{-1}v_-)$ on the fibers. Denote by $Z_m$
the space of compact effective one-cycles in $Y$ of degree $m$ over
$C$. Every irreducible component maps finitely onto $C$. Indeed, an
affine plane fiber contains no complete curve, and the projection
of each compact component is proper. The space $Z_0$ is the point
representing the empty cycle. Since $K_Y$ is trivial, all these
cycles have Chern number zero and belong to the semi-Fano locus.

\subsection{Symmetric functions on the fibers}
For $p\in C$ the fiber intersection of $D\in Z_m$ is a multiset of
$m$ points in $Y_p$. Project it to each of $L_{+,p}$ and $L_{-,p}$.
For $1\leq j\leq m$ the $j$th elementary symmetric function of a
projection gives an element
\[
 b_{\pm,j}(D)(p)\in L_{\pm,p}^{\otimes j}.
\]
These evaluations are continuous functions of $D$ by the
fiber-intersection map of \cite[(4.15)]{Pardon}.

\begin{lemma}\label{lem:support}
For $m>0$, choose $1+\max\{0,m(g-1)\}$ distinct points of $C$.
The common zero set of the functions $b_{\pm,j}$ evaluated at these
points consists of $m[C]$. This cycle is also the unique fixed point
of the circle on $Z_m$.
\end{lemma}
\begin{proof}
The coefficients $b_{\pm,j}(D)$ are holomorphic sections on $C$.
To see this across a branch point, normalize each irreducible
component. Its finite map to the smooth base curve is flat. In a
local trivialization of the line, the fiber coordinate acts by
multiplication on the finite locally free direct image of the
structure sheaf. Its characteristic polynomial records the projected
branches with multiplicities. Take the product of these polynomials
over the components, raising each to its cycle multiplicity. The
$j$th coefficient, up to the sign $(-1)^j$, transforms as a
section of $L_\pm^{\otimes j}$.

Both line bundles have degree $g-1$. If all the chosen evaluations
vanish, each coefficient vanishes identically by the degree bound.
The two projected polynomials are thus $T^m$ on every fiber. All
points of the multiset have both coordinates zero, which forces
$D=m[C]$. Conversely $m[C]$ is a common zero. For a circle-fixed
cycle the circle preserves each finite fiber multiset. Connectedness
forces it to fix each point individually. The origin is the only
fixed point of the fiber action.
\end{proof}

For later use, $Z_m$ is finite dimensional. Embed $Y$ in the bundle
of lines in $\cO_C\oplus L_+\oplus L_-$. On the open chart $Y$
the relative $\cO(1)$ is trivial. Tensoring it with a sufficiently
ample line bundle pulled back from $C$ gives a projective
polarization whose degree on $D$ is $m$ times its degree on the
zero section. Thus $Z_m$ is the open subset of a fixed-degree
projective Chow space parametrizing cycles disjoint from the divisor
at infinity. This description gives the compact-cycle topology on
that open set and the bounded dimension needed for equivariant
localization.

\subsection{The finite Borel construction}
Write $B_N=\PP^N$, $h=c_1(\cO_{B_N}(1))$, and
\[
 Z_{m,N}=\big((\CC^{N+1}\setminus\{0\})\times Z_m\big)/\CC^\times
        \longrightarrow B_N
\]
for the diagonal quotient. It is the relative degree-$m$ cycle
space of the corresponding Borel family of local threefolds.
Trivializing $L_\pm$ at each chosen point makes the tests of
Lemma~\ref{lem:support} into functions of weights $+j$ and $-j$.
A function $f$ of weight $r$ gives a section of the pullback of
$\cO_{B_N}(r)$. Our sign convention follows from the quotient:
the class of $(v,s)$ of weight one defines the functional on
$\CC v$ taking $v$ to $s$.

For any tuple $F=(f_1,\ldots,f_n)$ of equivariant functions of
nonzero weights $r_i$ with compact common zero set, put
\begin{equation}\label{eq:cut}
 u_N(F)=\frac{1}{\prod_i r_i}\,
    F^*\operatorname{Th}\!\left(\bigoplus_{i=1}^{n}\cO_{B_N}(r_i)\right)
        \quad\text{in } H_c^{2n}(Z_{m,N};\QQ).
\end{equation}
Here the Thom class is taken relative to the complement of the
zero section. Its pullback has support in the common zero set over
$B_N$, which is compact. In particular the definition allows
noncompact individual zero sets and needs no transversality. On
forgetting supports the normalized Euler class of a single section
is $c_1(\cO(r))/r=h$ for either sign of $r$.

\begin{lemma}\label{lem:stabilization}
The universal image of $u_N(F)$ depends only on $n-N$ for the given
local model and degree.
\end{lemma}
\begin{proof}
For the inclusion $i:B_N\hookrightarrow B_{N+1}$ of a hyperplane
use the same symbol for the inclusion of Borel families.
The class $u_{N+1}(F)$ restricts to $u_N(F)$, so the hyperplane
Gysin map gives
\[
 i_!u_N(F)=h\,u_{N+1}(F)=u_{N+1}(F,f)
\]
for every additional function $f$ of nonzero weight. The second
equality uses the compact support already supplied by $F$, which
allows the last relative Euler class to be replaced by its ordinary
class. The first equality is the projection formula for the complex
normal line pulled back from the base. The product description
locally over the base suffices even for a singular cycle space.

These Gysin maps induce the stabilization of relative classes in
\cite[(4.5)--(4.7)]{Pardon}. Given two tuples of lengths $n,n'$ at
stages $N,N'$ with $n-N=n'-N'$, append the second tuple to the
first and the first to the second. Both products now live at
stage $N+n'=N'+n$, and the even Thom factors commute. Hence the
images in the universal group coincide.
\end{proof}

Take $n=N$ in \eqref{eq:cut} and let $a_m(C,L)$ denote the image
in Pardon's algebra. Such a tuple exists by Lemma~\ref{lem:support}.
The oriented base has real dimension $2N$, so a compact class of
degree $2n$ has relative degree $2n-2N$ \cite[Lemma~3.5]{Pardon}.
The class $a_m(C,L)$ has virtual dimension and Chern number zero
and therefore lies in $H$. In particular its expression in the
scalar generators is a polynomial. The counit evaluates on the
empty cycle, hence vanishes on $a_m(C,L)$ for $m>0$. We set
$a_0(C,L)=1$, as given by the empty tuple at $N=0$.

\section{Addition and multiplication of cycles}\label{sec:operations}

The coproduct $\Delta$ in $H$ is defined by pullback under cycle
addition; the operation $\rho_d$ is pullback under $D\mapsto dD$,
for $d\geq1$ \cite[Definitions~3.9 and~3.12]{Pardon}. The two
arguments of the coproduct are cycles on independent families of
threefolds. In the Borel construction this distinction introduces
the class of the diagonal in a product of two projective spaces.

\begin{proposition}\label{prop:operations}
For a fixed balanced pair $(C,L)$, write $a_m=a_m(C,L)$. Then
\begin{equation}\label{eq:operations}
 \Delta a_m=\sum_{s+t=m}a_s\otimes a_t,\qquad
 \rho_d a_m=
 \begin{cases}a_{m/d}&d\mid m,\\0&d\nmid m.\end{cases}
\end{equation}
\end{proposition}
\begin{proof}
The preimage of $Z_m$ under multiplication by $d$ is empty unless
$d\mid m$. Otherwise pull the tuple of Lemma~\ref{lem:support}
back to $Z_{m/d}$. The pulled-back functions have the same weights,
with common zero $(m/d)[C]$. Naturality of the relative Thom class
and Lemma~\ref{lem:stabilization} give the formula for $\rho_d$.

For the coproduct fix $s+t=m$. Over the diagonal in $B_N^2$, addition
gives a map from the common-base family of $Z_s\times Z_t$ to
$Z_{m,N}$. In the two-base description of $\Delta$
\cite[(4.17)--(4.21)]{Pardon}, the $(s,t)$ component is represented
by pulling back $u_N(F)$ along addition and then applying the
oriented Gysin map $\delta_{N!}$ over this diagonal. Its target is
the independent product $Z_{s,N}\times Z_{t,N}$. We can take
$|F|=N$ with common zero $m[C]$. The tuple pulled back along
addition has common zero $(s[C],t[C])$, since both summands are
effective cycles.

Choose compactly supported cut-downs given by tuples $F_s,F_t$ on
the two factors, of lengths $n_s,n_t$, and set $M=n_s+n_t$. On
the common-base family their concatenation also has compact common
zero set. The proof of Lemma~\ref{lem:stabilization}, applied to
the diagonal action on $Z_s\times Z_t$, identifies its class at
stage $M$ with the pulled-back class at stage $N$ after
stabilization. We may make this substitution before applying the
diagonal Gysin map. Indeed, for the simultaneous hyperplane
inclusion $i$ the identity
\[
 (i\times i)_!\,\delta_{N!}=\delta_{N+1,!}\,i_!
\]
follows by composition of oriented embeddings on the bases.
Hyperplane pushforward in each of the two independent bases
preserves the image of the corresponding relative class.

At stage $M$ the concatenated class on the diagonal is the
restriction of $u_M(F_s)\boxtimes u_M(F_t)$. By the projection
formula its diagonal pushforward is this external product times
\[
 [\delta_M(B_M)]=\sum_{i+j=M}h_1^ih_2^j .
\]
The two factors of the $(i,j)$ summand have relative cohomological
degrees $2(n_s+i-M)$ and $2(n_t+j-M)$, whose sum is zero.
A positive degree here is a negative virtual dimension. Its image
in the semi-Fano algebra vanishes by \cite[Theorem~1.2]{Pardon},
as in the proof of \cite[Theorem~7.6]{Pardon}. Consequently only
$i=n_t$, $j=n_s$ survives. Multiplication by these powers of the
hyperplane classes completes each tuple to length $M$ at stage
$M$, giving $a_s\otimes a_t$ with coefficient one.
An empty factor admits the empty tuple and the same argument, with
any extra copies of $h$ represented by the Euler class of
$\cO(1)$. Summation over $s+t=m$ proves \eqref{eq:operations}.
\end{proof}

\section{Local Gromov--Witten evaluation and reconstruction}\label{sec:evaluation}

Let $u$ be the Gromov--Witten variable and $q$ the local curve-degree
variable. We use disconnected maps with no wholly contracted
connected components. Contracted components attached to positive-degree
maps are included. For each partition $\lambda$ let $h(b)$ denote the
hook length of a box $b$ of its Young diagram. The empty partition
contributes $1$ to the following series:
\begin{equation}\label{eq:hook}
 E_g(q,u)=\sum_{\lambda} q^{|\lambda|}
       \prod_{b\in\lambda}\bigl(2\sin(h(b)u/2)\bigr)^{2g-2}.
\end{equation}

\begin{proposition}\label{prop:evaluation}
The disconnected Gromov--Witten evaluation of
$\sum_{m\ge0}a_m(C,L)q^m$ is $E_g(q,u)$.
\end{proposition}
\begin{proof}
Fix $m>0$ and the Euler characteristic of the source.
The stable-map cycle map is proper at fixed Euler characteristic,
and its equivariant virtual pushforward defines the Gromov--Witten
class on $Z_m$ \cite[\S3.5, (3.21)--(3.23)]{Pardon}.
By Lemma~\ref{lem:support} its fixed-locus localization is a class
$v$ at $m[C]$. More precisely, equivariant homology after inverting
$h$ identifies with that of the fixed locus by
\cite[Proposition~4.5]{Pardon}, and $v$ is the inverse image under
fixed-locus pushforward. Indeed, off $m[C]$ the tuple of
Lemma~\ref{lem:support} is a nowhere-zero section. Its Euler
class $(\prod_i r_i)h^n$ vanishes there, so $h^n$ annihilates
the equivariant homology of the complement. The fixed-locus
isomorphism follows from the localization sequence. The class $v$
is a scalar because the virtual dimension is zero.

We check that the compact class \eqref{eq:cut} makes precisely this
scalar an ordinary universal evaluation. Fix a tuple of length $n$
as in Lemma~\ref{lem:support}. At each Borel stage its compact class
restricts on the fixed section to $h^n$. Thus capping with the
equivariant virtual class and pushing forward to the point gives
$h^n v$ after inverting $h$, by the projection formula. The compact
support of the cut-down already defines this pushforward over
$\QQ[h]$, which injects into $\QQ[h,h^{-1}]$. At stage $N=n$ the
ordinary integral is $v$ since $\int_{B_n}h^n=1$. This is the
evaluation of $a_m(C,L)$, by the finite Borel description in
\cite[Proposition~4.8]{Pardon}.

It remains to identify $v$. A fixed map has fixed image cycle, so
all its positive-degree components map into $C$. Attached contracted
components map into $C$ as well. Conversely every stable map into
$C$ is fixed. Write $\mathcal F$ for this proper fixed moduli stack
and $p$ for the cycle map. Along $C$ the tangent bundle splits as
$T_C\oplus L_+\oplus L_-$. The fixed part of the obstruction theory
is the stable-map theory of $C$; the moving virtual normal complex
is $R\pi_*f^*L_+\oplus R\pi_*f^*L_-$ with the respective
characters $+1,-1$, where $(\pi,f)$ is the universal curve and map.
Both characters stay nonzero on the antidiagonal. Equivariant
virtual localization and proper pushforward give
\[
 v=(p|_{\mathcal F})_*
    \left(\frac{[\mathcal F]^{\mathrm{vir}}}
    {e_{S^1}(N^{\mathrm{vir}})}\right).
\]
This is the residue integral of \cite[\S2.2]{BP}. Its degree-$m$
normalization shifts the power of $u$ by
$m(2-2g+\deg L_++\deg L_-)=0$. Formula~\eqref{eq:hook} now
follows from \cite[Corollary~7.3]{BP}.
\end{proof}

The degree-two term distinguishes the balanced degrees among all
splittings with determinant $K_C$. Indeed, take line degrees
$g-1+\ell,g-1-\ell$ and retain the antidiagonal action. Write
$E_{g,\ell}$ for its disconnected residue series, so $E_{g,0}=E_g$.
In the formula of \cite[Corollary~7.2]{BP} the change of degrees in
degree two multiplies the two partition contributions by
$e^{i\ell u}$ and $e^{-i\ell u}$. Indeed, the partitions $(2)$ and
$(1,1)$ have opposite total contents $\pm1$ and the same hook
lengths, and the sign in degree two is $(-1)^{2\ell}=1$. Thus
\begin{equation}\label{eq:balance}
 [q^2]E_{g,\ell}=\cos(\ell u)\,[q^2]E_g .
\end{equation}
The balanced coefficient is a nonzero Laurent series. Equality
with the elementary-cluster coefficient therefore forces $\ell=0$.

To pass from this evaluation to a class in $H$, we extend the
reconstruction argument of \cite[\S2]{Transition}. Write
$x_{h,j}=x_{h,j,0}$ for Pardon's scalar generators,
$x_{h,0}=1$, and set
\[
 f_h(q)=\sum_{j\geq0}x_{h,j}q^j,\qquad
 p_{h,j}=[q^j]\log f_h(q)\quad(j>0).
\]
Over $\QQ$ the $p_{h,j}$ are polynomial coordinates on $H$.
Pardon's scalar formulas \cite[Lemmas~4.12--4.13 and
Theorem~7.6]{Pardon} give \eqref{eq:operations} for $x_{h,m}$.
Thus $p_{h,j}$ is primitive for $\Delta$, and $\rho_d p_{h,j}$
is $p_{h,j/d}$ if $d\mid j$ and zero otherwise.

Let $\Gamma$ be a submonoid of a finitely generated abelian group with
an additive degree on $\Gamma$, positive and integral off zero, and
with finite sublevel sets. Complete the monoid ring in this degree.
Call a series
$U=\sum_\beta U_\beta Q^\beta$ \emph{compatible} if $U_0=1$, its
other coefficients lie in the augmentation ideal of $H$, and
\begin{equation}\label{eq:compat}
 \Delta U_\beta=\sum_{\gamma+\eta=\beta}U_\gamma\otimes U_\eta,
 \qquad \rho_d U_\beta=\sum_{\substack{\gamma\in\Gamma\\ d\gamma=\beta}}U_\gamma .
\end{equation}
The second identity says that $\rho_d U=\psi_d U$, where
$\rho_d$ acts on coefficients and $\psi_d$ substitutes
$Q^\gamma\mapsto Q^{d\gamma}$. The latter is a continuous ring
endomorphism even when multiplication by $d$ identifies distinct
classes. The coefficient at $\beta$ sums all their contributions.

\begin{lemma}\label{lem:reconstruction}
A compatible series is determined by its disconnected
Gromov--Witten evaluation. It admits a unique product expansion
in the $f_h(Q^\beta)$ with rational exponents and finite genus
support in each class.
\end{lemma}
\begin{proof}
The primitive polynomials in the $p_{h,j}$ are precisely their
finite linear combinations, since an additive polynomial over
$\QQ$ is linear. If a primitive $R$ is killed by every $\rho_d$
for $d>1$, it is a linear combination of $p_{h,1}=x_{h,1}$.
Indeed applying $\rho_j$ for the largest multiplicity $j>1$ in
such a polynomial detects all its coefficients at $j$. Set
$w=(2\sin(u/2))^2$. The scalar degree-one evaluations are
$\GW(x_{h,1})=w^{h-1}$ by the degree-one local formula
\cite[\S8]{BP}. These functions are linearly independent by their
distinct lowest powers of $u$. If two compatible series differ,
choose a class $\beta$ of least degree at which they differ.
Every decomposition of $\beta$ into two nonzero summands uses
lower degrees. The first identity in \eqref{eq:compat} therefore
makes the difference at $\beta$ primitive. For $d>1$ each root
$d\gamma=\beta$ has degree $\deg(\beta)/d$, so the second identity
kills the difference under $\rho_d$. Equal evaluations force it
to vanish.

For the expansion, every coefficient of $\log U$ is primitive.
Indeed the coproduct identity gives $\Delta U=(U\otimes1)(1\otimes U)$
with one common monoid variable, so taking the logarithm gives the
primitive identity coefficientwise. Write
$[Q^\beta]\log U=\sum_{h,j}b_{\beta;h,j}p_{h,j}$. Since $\rho_j$
and $\psi_j$ commute with logarithms, comparison of the $p_{h,1}$
coefficients gives
\[
 b_{\beta;h,j}=\sum_{j\gamma=\beta} b_{\gamma;h,1}.
\]
This is exactly the $p_{h,j}Q^\beta$ coefficient of
$\sum_{\gamma\ne0,h}b_{\gamma;h,1}\log f_h(Q^\gamma)$.
The sums are finite in each degree, because positive-class degrees
are at least one and each coefficient in $H$ is a polynomial.
Exponentiating gives the expansion. Its exponents are unique since
$p_{h,1}Q^\gamma$ in the logarithm recovers $b_{\gamma;h,1}$.
\end{proof}

\begin{corollary}\label{cor:characterization}
Assume the universal Gromov--Witten evaluation for analytic families.
For each $g$ there is exactly one compatible series in $H\llbracket q\rrbracket$
whose disconnected evaluation is the contribution of a genus-$g$
elementary cluster. Every balanced antidiagonal model of genus $g$
represents this series.
\end{corollary}
\begin{proof}
Propositions~\ref{prop:operations} and~\ref{prop:evaluation} give
existence, since the elementary-cluster contribution is $E_g$
\cite[Proposition~3.3]{IP}. Lemma~\ref{lem:reconstruction} with
$\Gamma=\mathbb N$ gives uniqueness. In particular the coefficients
are independent of the curve, the line bundle, and the order of
the two summands. We denote them by $a_{g,m}$ and the series by $A_g$.
\end{proof}

\section{The product expansion on a compact threefold}\label{sec:factor}

Let $X$ be a smooth projective Calabi--Yau threefold. Choose a
projective embedding of $X$ and use its degree on $\Gamma_X$.
Every nonempty effective cycle has positive degree. The integral
homology class is locally constant on the cycle space
\cite[\S2.1]{Pardon}. The Chow variety of cycles of each fixed degree
is compact, so its image in the discrete group $H_2(X;\ZZ)$ is finite.
Thus $\Gamma_X$ has finite degree sublevel sets, and the locus
$\mathscr C_\beta(X)$ of cycles of class $\beta$ is compact and
open and closed, including when different classes have the same
image modulo torsion.

Put $U_{X,\beta}=(X,1_{\mathscr C_\beta(X)})\in H$. Addition pulls the
characteristic function back to the sum over decompositions of $\beta$.
Under multiplication $D\mapsto dD$, the inverse image of
$\mathscr C_\beta(X)$ is the disjoint union of the
$\mathscr C_\gamma(X)$ over all effective $\gamma$ with $d\gamma=\beta$.
Pullback of its characteristic function counts each locus once.
This proves \eqref{eq:compat} for $\Z_X=\sum_\beta U_{X,\beta}Q^\beta$.

\begin{proof}[Proof of Theorem~\ref{thm:main}]
Use the symplectic form of the polarization. Proposition~3.3 and
Theorem~8.1 of Ionel and Parker \cite{IP} and the projective
comparison theorem of Siebert \cite[Theorem~0.1]{Siebert} give
\begin{equation}\label{eq:cluster}
 \log\GW(\Z_X)=
    \sum_{\beta\ne0}\sum_{g\ge0}c_{\beta,g}(X)\log E_g(Q^\beta,u).
\end{equation}
Here $\log\GW(\Z_X)$ is the connected positive-class series. The
integers in \eqref{eq:cluster} have finite genus support at each
$\beta$ by \cite[Definition~3.45 and Theorem~3.50]{DIW},
which gives finiteness for this elementary expansion itself.
Both the cluster expansion and the comparison of virtual classes
are indexed by $H_2(X;\ZZ)$. In particular \cite[(8.9)]{IP} sums
over every pair $(d,B)$ with $dB=A$ in integral homology. Siebert's
comparison fixes the integral class and the connected genus. The
semipositive description of the symplectic invariants is recalled
in \cite[Remark~1.1]{IP}.

Only classes of positive symplectic area can occur in the sum.
Suppose $\beta$ has positive ample degree and is ineffective for
the projective structure on $X$. Every $\gamma$ with $d\gamma=\beta$,
$d>1$, is then ineffective, since a positive multiple of an
effective cycle is effective. There are finitely many such pairs
because $d$ divides $\deg(\beta)$ and the roots for a fixed $d$,
if any, form a coset of the finite kernel of $d$ on $H_2(X;\ZZ)$.
Compare coefficients in
\eqref{eq:cluster}, inductively on the positive integral degree.
The left side vanishes at $\beta$, and all terms from proper
divisors already vanish. Thus
$\sum_g c_{\beta,g}(X) w^{g-1}=0$, using
$[q]\log E_g(q,u)=w^{g-1}$. Linear independence gives
$c_{\beta,g}(X)=0$ for every $g$.

The right side of \eqref{eq:main} is now well defined in
$H\llbracket\Gamma_X\rrbracket$. There are finitely many effective
classes of bounded degree, finitely many contributing genera per
class, and polynomial coefficients in each factor. Products and
inverses preserve the coproduct identity. By
\eqref{eq:operations}, $\rho_d A_g(Q^\beta)=A_g(Q^{d\beta})
=\psi_d A_g(Q^\beta)$. Both $\rho_d$ and $\psi_d$ are ring
endomorphisms, so the full product also satisfies the root sum
in \eqref{eq:compat}. Proposition~\ref{prop:evaluation} and
\eqref{eq:cluster} identify the disconnected GW evaluation of
the product with $\GW(\Z_X)$. Lemma~\ref{lem:reconstruction}
identifies the universal series. Uniqueness follows upon evaluation
from the unique elementary expansion, or by comparing the
multiplicity-one term at a first class of difference.
\end{proof}

The elementary-cluster evaluation in Corollary~\ref{cor:characterization}
is essential to the characterization. For example, fix $g\ne h$ and
replace $A_g(q)$ by $A_g(q)A_h(q^2)$, keeping every other factor.
The new series is compatible and has the same multiplicity-one
coefficient as $A_g$. In the compact product its exponent stays
$c_{\beta,g}$, and the exponent of $A_h(Q^\beta)$ becomes
$c_{\beta,h}-\sum_{2\gamma=\beta}c_{\gamma,g}$. Thus this change
preserves integer exponents in every class, although it changes the
degree-two evaluation of the genus-$g$ factor.

\subsection{Integrality from stable pairs}
There is a second way to see why the choice of action in
\eqref{eq:main} yields integer exponents. It uses the
Gromov--Witten/stable-pair correspondence for $H$ and the
form of the local hook series. We first obtain rational exponents
with finite genus support without using the numerical
cluster decomposition.

By Lemma~\ref{lem:reconstruction} the compatible $A_g(q)$ is a
product of rational powers of $f_h(q^d)$. Its degree-one evaluation
is $w^{g-1}$. Thus the entire factor at $d=1$ is $f_g(q)$, by the
independence of the powers $w^{h-1}$. All other factors have
$d\geq2$. At each multiplicity $m$ this gives
\begin{equation}\label{eq:triangular}
 a_{g,m}=x_{g,m}+P_{g,m}\bigl(\{x_{h,j}:1\le j<m\}\bigr)
\end{equation}
with $P_{g,m}$ a finite rational polynomial. Induction on $m$
inverts the substitution $x_{g,m}\mapsto a_{g,m}$. It is therefore
an automorphism of the polynomial algebra $H$. The identities
\eqref{eq:operations} show that it intertwines both the coproduct
and each $\rho_d$. Apply Lemma~\ref{lem:reconstruction} to the
inverse image of $\Z_X$, then apply the automorphism. This gives
a product in the $A_g(Q^\beta)$ with rational exponents
$n_{\beta,g}$ and finite genus support in every class.

Write $z$ for the stable-pair variable. Pardon's Theorem~1.7
\cite{Pardon}, extended by linearity over $\QQ$, identifies
$\PT(a_{g,m})$ with the expansion at $z=0$ of the rational
function $[q^m]E_g(q,u)$ under $e^{iu}=-z$. The identification
is unambiguous since expansion at $e^{iu}=1$ determines a rational
function. For a box of hook length $k$ one has
\[
 (2\sin(ku/2))^2=2-(-z)^k-(-z)^{-k}
                 =-(-z)^{-k}(1-(-z)^k)^2 .
\]
This is a unit in $\ZZ((z))$. Hence every coefficient of
$\PT(A_g(q))$ lies in $\ZZ((z))$, including when $g=0$ and
the hook factors are inverted. The compact stable-pair invariants
$\PT(U_{X,\beta})$ also lie in $\ZZ((z))$
\cite[\S3.5]{Pardon}.

Suppose all exponents in classes of smaller ample degree than
$\beta$ are known to be integral. Divide the stable-pair series by
the corresponding powers of the local factors. Their constant
terms are one, so the residual series still has coefficients in
$\ZZ((z))$. Its coefficient at $\beta$ is
\begin{equation}\label{eq:residual}
 \sum_g n_{\beta,g}\,z^{1-g}(1+z)^{2g-2},
\end{equation}
by the degree-one hook formula. In this finite sum the highest
genus with nonzero exponent has the unique lowest Laurent power,
and its local factor has leading coefficient one. The coefficient
of that power is therefore the exponent itself. It is integral.
Subtract its summand in \eqref{eq:residual} and continue in
descending genus. This proves integrality for every $\beta$.
Equation~\eqref{eq:cluster} identifies these integers as
$c_{\beta,g}(X)$.

\section{Comparison with scalar local curves}\label{sec:comparison}

The first discrepancy between the two actions can be read in
$H$ itself. Equation~\eqref{eq:triangular}, or reconstruction in
degree one, gives $a_{g,1}=x_{g,1}$. Thus the difference
$a_{2,2}-x_{2,2}$ is primitive: the mixed coproduct term in
each case is $x_{2,1}\otimes x_{2,1}$. It is killed by all
$\rho_d$ with $d>1$ and hence is a finite linear combination of
$x_{h,1}$. Each partition of two has hook lengths $1,2$. Its
genus-two contribution to \eqref{eq:hook} is $w^2(4-w)$.
At scalar weights the corresponding total is $8w^2$
\cite[\S8]{BP}. The difference of the evaluations is $-2w^3$,
which is the evaluation of $-2x_{4,1}$. It follows that
\begin{equation}\label{eq:two}
 a_{2,2}=x_{2,2}-2x_{4,1}.
\end{equation}

In degree three the change of local factors affects the integrality
of compact product exponents. We give the comparison to fix the
direction of the change. Let $e_{\beta,g}$ be the exponent of
$f_g(Q^\beta)$ in the scalar factorization of $\Z_X$, and put
\[
 F_g=\GW(f_g),\qquad K_g=\log F_g,\qquad J_g=\log E_g .
\]
At an effective class $\alpha$ primitive in $H_2(X;\ZZ)$, the
multiplicity-one terms give $e_{\alpha,g}=c_{\alpha,g}$. The only
proper integral divisor of $3\alpha$ is $\alpha$. Comparison of
the two connected series at this class gives
\[
 \sum_h(c_{3\alpha,h}-e_{3\alpha,h})w^{h-1}
       =\sum_g c_{\alpha,g}[q^3](K_g-J_g).
\]
The scalar local calculation of \cite[\S3]{Counterexample} gives
\begin{align*}
 [q^3](K_2-J_2)&=120w^4-20w^5+2w^6,\\
 [q^3](K_3-J_3)&=14448w^7-1136w^8+\tfrac{4720}{3}w^9
                       -332w^{10}+40w^{11}-2w^{12}.
\end{align*}
For the second row the scalar disconnected coefficient is
$2673w^6+9612w^7+2534w^8+(292/3)w^9+w^{10}$. The three
antidiagonal partitions contribute $2U^2+V^2$, with
$U=w^3(4-w)(3-w)^2$ and $V=w^3(3-w)^2$.
Passing to the logarithmic difference subtracts the quadratic
cross term $w^2(24w^5-2w^6)$. For $g=0,1$ the local series
agree, and for $g\ge4$ their cubic difference is
$O(u^{6(g-1)+2})$ \cite[\S2]{Counterexample}. Thus only the
genus-three row contributes to $w^9$. This recovers
\begin{equation}\label{eq:three}
 e_{3\alpha,10}=c_{3\alpha,10}-\tfrac{4720}{3}c_{\alpha,3}.
\end{equation}

On the smooth projective toric hypersurface of \cite{Counterexample}
there is a first Hirzebruch surface with negative section $s$ and
ruling fiber $f$. The class $\alpha=3s+4f$ is primitive. The
minimal stable-pair moduli space in this class is $\PP^{13}$,
and the primitive genus-three coefficient is $c_{\alpha,3}=-14$.
For this threefold \eqref{eq:three} becomes
\[
 e_{3\alpha,10}=c_{3\alpha,10}+\frac{66080}{3}.
\]
The antidiagonal exponent on the right is an integer by
Theorem~\ref{thm:main}, while the scalar exponent on the left
has fractional part $2/3$. Both products represent the same
universal cycle series.

\begin{thebibliography}{DIW26}
\bibitem[BP08]{BP} J.~Bryan and R.~Pandharipande,
\emph{The local Gromov--Witten theory of curves},
J. Amer. Math. Soc. \textbf{21} (2008), 101--136,
\href{https://arxiv.org/abs/math/0411037v3}{arXiv:math/0411037v3}.
\bibitem[DIW26]{DIW} A.~Doan, E.-N.~Ionel and T.~Walpuski,
\emph{The Gopakumar--Vafa finiteness conjecture},
Ann. of Math. (2) \textbf{204} (2026), 819--871,
\href{https://doi.org/10.4007/annals.2026.204.3.1}{doi:10.4007/annals.2026.204.3.1}.
Theorem numbering follows \href{https://arxiv.org/abs/2103.08221v2}{arXiv:2103.08221v2}.
\bibitem[IP18]{IP} E.-N.~Ionel and T.~H.~Parker,
\emph{The Gopakumar--Vafa formula for symplectic manifolds},
Ann. of Math. (2) \textbf{187} (2018), 1--64,
\href{https://arxiv.org/abs/1306.1516v6}{arXiv:1306.1516v6}.
\bibitem[Par25]{Pardon} J.~Pardon,
\emph{Universally counting curves in Calabi--Yau threefolds},
\href{https://arxiv.org/abs/2308.02948v3}{arXiv:2308.02948v3} (2025).
\bibitem[Sie99]{Siebert} B.~Siebert,
\emph{Algebraic and symplectic Gromov--Witten invariants coincide},
Ann. Inst. Fourier (Grenoble) \textbf{49} (1999), 1743--1795,
\href{https://arxiv.org/abs/math/9804108v1}{arXiv:math/9804108v1}.
\bibitem[Wue26a]{Counterexample} B.~J.~Wuebben,
\emph{A counterexample to integral product exponents in universal curve counting},
preprint (2026), \href{https://doi.org/10.5281/zenodo.23750048}{doi:10.5281/\allowbreak zenodo.23750048}.
\bibitem[Wue26b]{Transition} B.~J.~Wuebben,
\emph{Universal curve counting under conifold transitions},
preprint (2026), \href{https://doi.org/10.5281/zenodo.23750046}{doi:10.5281/\allowbreak zenodo.23750046}.
\end{thebibliography}

\end{document}
